\documentclass[10pt,a4paper]{amsart}
\usepackage{amsmath,amssymb,amsthm,graphicx}
\usepackage{verbatim}
\newenvironment{reference}{\comment}{\endcomment}
%\newenvironment{reference}{}{}
\newenvironment{slogan}{\comment}{\endcomment}
\newenvironment{history}{\comment}{\endcomment}

% For commutative diagrams we use Xy-pic
\usepackage[all]{xy}

% We use 2cell for 2-commutative diagrams.
\xyoption{2cell}
% Theorem environments.
%
\theoremstyle{plain}
\newtheorem{theorem}[subsection]{Theorem}
\newtheorem{proposition}[subsection]{Proposition}
\newtheorem{lemma}[subsection]{Lemma}

\theoremstyle{definition}
\newtheorem{definition}[subsection]{Definition}
\newtheorem{example}[subsection]{Example}
\newtheorem{exercise}[subsection]{Exercise}
\newtheorem{situation}[subsection]{Situation}

\theoremstyle{remark}
\newtheorem{remark}[subsection]{Remark}
\newtheorem{remarks}[subsection]{Remarks}

\numberwithin{equation}{subsection}

% Macros
%
\def\lim{\mathop{\mathrm{lim}}\nolimits}
\def\colim{\mathop{\mathrm{colim}}\nolimits}
\def\Spec{\mathop{\mathrm{Spec}}}
\def\Hom{\mathop{\mathrm{Hom}}\nolimits}
\def\Ext{\mathop{\mathrm{Ext}}\nolimits}
\def\SheafHom{\mathop{\mathcal{H}\!\mathit{om}}\nolimits}
\def\SheafExt{\mathop{\mathcal{E}\!\mathit{xt}}\nolimits}
\def\Sch{\mathit{Sch}}
\def\Mor{\mathop{\mathrm{Mor}}\nolimits}
\def\Ob{\mathop{\mathrm{Ob}}\nolimits}
\def\Sh{\mathop{\mathit{Sh}}\nolimits}
\def\NL{\mathop{N\!L}\nolimits}
\def\CH{\mathop{\mathrm{CH}}\nolimits}
\def\proetale{{pro\text{-}\acute{e}tale}}
\def\etale{{\acute{e}tale}}
\def\QCoh{\mathit{QCoh}}
\def\Ker{\mathop{\mathrm{Ker}}}
\def\Im{\mathop{\mathrm{Im}}}
\def\Coker{\mathop{\mathrm{Coker}}}
\def\Coim{\mathop{\mathrm{Coim}}}

% Boxtimes
%
\DeclareMathSymbol{\boxtimes}{\mathbin}{AMSa}{"02}

%
% Macros for moduli stacks/spaces
%
\def\QCohstack{\mathcal{QC}\!\mathit{oh}}
\def\Cohstack{\mathcal{C}\!\mathit{oh}}
\def\Spacesstack{\mathcal{S}\!\mathit{paces}}
\def\Quotfunctor{\mathrm{Quot}}
\def\Hilbfunctor{\mathrm{Hilb}}
\def\Curvesstack{\mathcal{C}\!\mathit{urves}}
\def\Polarizedstack{\mathcal{P}\!\mathit{olarized}}
\def\Complexesstack{\mathcal{C}\!\mathit{omplexes}}
% \Pic is the operator that assigns to X its picard group, usage \Pic(X)
% \Picardstack_{X/B} denotes the Picard stack of X over B
% \Picardfunctor_{X/B} denotes the Picard functor of X over B
\def\Pic{\mathop{\mathrm{Pic}}\nolimits}
\def\Picardstack{\mathcal{P}\!\mathit{ic}}
\def\Picardfunctor{\mathrm{Pic}}
\def\Deformationcategory{\mathcal{D}\!\mathit{ef}}
\usepackage{cleveref}
\crefname{originalreference}{original reference}{original references}
\Crefname{originalreference}{Original reference}{Original references}
\makeatletter
\newlabel{lemma-etale-image-open}{{06NF}{source}{}{sourcepointer-1}{}}
\newlabel{lemma-etale-image-open@cref}{{[originalreference][0][]06NF}{[1][0][]source}{}{}{}}
\makeatother
\crefname{originallemma}{original lemma}{original lemmas}
\Crefname{originallemma}{Original lemma}{Original lemmas}
\makeatletter
\newlabel{schemes-lemma-compose-after-separated}{{01KV}{source}{}{sourcepointer-2}{}}
\newlabel{schemes-lemma-compose-after-separated@cref}{{[originallemma][0][]01KV}{[1][0][]source}{}{}{}}
\makeatother
\crefname{originallemma}{original lemma}{original lemmas}
\Crefname{originallemma}{Original lemma}{Original lemmas}
\makeatletter
\newlabel{schemes-lemma-immersions-monomorphisms}{{01L7}{source}{}{sourcepointer-3}{}}
\newlabel{schemes-lemma-immersions-monomorphisms@cref}{{[originallemma][0][]01L7}{[1][0][]source}{}{}{}}
\makeatother
\crefname{originallemma}{original lemma}{original lemmas}
\Crefname{originallemma}{Original lemma}{Original lemmas}
\makeatletter
\newlabel{schemes-lemma-separated-permanence}{{01KU}{source}{}{sourcepointer-4}{}}
\newlabel{schemes-lemma-separated-permanence@cref}{{[originallemma][0][]01KU}{[1][0][]source}{}{}{}}
\makeatother
\crefname{originallemma}{original lemma}{original lemmas}
\Crefname{originallemma}{Original lemma}{Original lemmas}
\makeatletter
\newlabel{spaces-lemma-composition-representable-transformations-property}{{02WK}{source}{}{sourcepointer-5}{}}
\newlabel{spaces-lemma-composition-representable-transformations-property@cref}{{[originallemma][0][]02WK}{[1][0][]source}{}{}{}}
\makeatother
\crefname{originallemma}{original lemma}{original lemmas}
\Crefname{originallemma}{Original lemma}{Original lemmas}
\makeatletter
\newlabel{spaces-lemma-properties-diagonal}{{02X4}{source}{}{sourcepointer-6}{}}
\newlabel{spaces-lemma-properties-diagonal@cref}{{[originallemma][0][]02X4}{[1][0][]source}{}{}{}}
\makeatother
\crefname{originallemma}{original lemma}{original lemmas}
\Crefname{originallemma}{Original lemma}{Original lemmas}
\makeatletter
\newlabel{spaces-lemma-representable-morphisms-spaces-property}{{02WZ}{source}{}{sourcepointer-7}{}}
\newlabel{spaces-lemma-representable-morphisms-spaces-property@cref}{{[originallemma][0][]02WZ}{[1][0][]source}{}{}{}}
\makeatother
\crefname{originallemma}{original lemma}{original lemmas}
\Crefname{originallemma}{Original lemma}{Original lemmas}
\makeatletter
\newlabel{spaces-lemma-representable-transformations-property-implication}{{02YO}{source}{}{sourcepointer-8}{}}
\newlabel{spaces-lemma-representable-transformations-property-implication@cref}{{[originallemma][0][]02YO}{[1][0][]source}{}{}{}}
\makeatother
\title{An etale chart by a separated scheme is a separated morphism and its fibre product is a separated scheme}
\author{The Stacks Project Authors}
\date{stacks-project:spaces-properties.tex:local-snapshot-20261012; GNU FDL 1.2 or later}
\begin{document}
\maketitle
\noindent\textit{Editorial note.} All delivered mathematical text is the author's own, copied from the pinned Stacks Project chapter file spaces-properties.tex; no statement or proof step was rewritten, reordered or removed.\par\medskip
\noindent\textit{Editorial note.} The counted proof is the author's complete argument: the reduction to a surjective chart using the separatedness of the open immersion onto the image, the proof that the fibre product is a separated scheme through the separatedness of the diagonal morphism and of the product over the base, and the closing application of the chart criterion.\par\medskip
\noindent\textit{Editorial note.} The label lemma-separated-cover names the counted statement, and the proof environment immediately following it in the source is delivered verbatim as its proof.\par\medskip
\noindent\textit{Editorial note.} The counted statement is carried by the author's own lemma environment in the chapter source.\par\medskip
\noindent\textit{Editorial note.} Pre-open gate: python operations/agent-authoring/tools/precard\_gate.py --source spaces-properties.tex --label lemma-separated-cover --start 991 --end 1036 --local-id 28051 returned RESULT=OK (L28051 free, occupancy FREE, public interval clear). Candidate taken from sources/stacks-spaces-properties/candidates-halfA.json.\par\medskip
\noindent\textit{Editorial note.} The statement and proof are copied from the pinned chapter file; the cited inputs are kept as external pointers with their canonical Stacks tags.\par\medskip
\noindent\textit{Allowed standard prerequisites (editorial).}\begin{itemize}
\item Factorization of an etale morphism through the open image. Same chapter, the lemma that an etale morphism from a scheme factors through a surjective etale morphism onto the open subspace which is its image. Supplies the reduction to the case of a surjective etale chart, the open immersion being separated.
\item Separation of the diagonal morphism of the relation. Spaces, the lemma on properties of the diagonal, whose proof shows that the map from the fibre product to the product over the base is separated. Supplies the separatedness of the morphism j : R -> U x\_S U.
\item Permanence of separatedness under composition and base change, and the criterion on a chart. Schemes, the lemmas that separatedness is stable under composition and base change, and Spaces, the lemma reducing properties of representable morphisms to a surjective etale chart. Supply the separatedness of U x\_S U and of R, and the conclusion that the chart U -> X is separated.
\end{itemize}\medskip
\ifdefined\hypertarget\hypertarget{sourcepointer-1}{}\fi
\noindent\textit{Source reference.} Original reference 06NF: sources/stacks-spaces-properties/pinned-original-source.tex, source lines 753--753. Prerequisite pointer invoked at proof line(s) 1002; it names a standard fact of the same chapter that the counted argument uses, and every step of the counted proof built on it is retained verbatim in the delivered body.\par\medskip
\ifdefined\hypertarget\hypertarget{sourcepointer-2}{}\fi
\noindent\textit{Source reference.} Original lemma 01KV: sources/stacks-chow/tags.txt, source lines 2064--2064. Prerequisite pointer invoked at proof line(s) 1026; it names a standard fact of a sibling chapter that the counted argument uses, and every step of the counted proof built on it is retained verbatim in the delivered body.\par\medskip
\ifdefined\hypertarget\hypertarget{sourcepointer-3}{}\fi
\noindent\textit{Source reference.} Original lemma 01L7: sources/stacks-chow/tags.txt, source lines 2076--2076. Prerequisite pointer invoked at proof line(s) 1010; it names a standard fact of a sibling chapter that the counted argument uses, and every step of the counted proof built on it is retained verbatim in the delivered body.\par\medskip
\ifdefined\hypertarget\hypertarget{sourcepointer-4}{}\fi
\noindent\textit{Source reference.} Original lemma 01KU: sources/stacks-chow/tags.txt, source lines 2063--2063. Prerequisite pointer invoked at proof line(s) 1028; it names a standard fact of a sibling chapter that the counted argument uses, and every step of the counted proof built on it is retained verbatim in the delivered body.\par\medskip
\ifdefined\hypertarget\hypertarget{sourcepointer-5}{}\fi
\noindent\textit{Source reference.} Original lemma 02WK: sources/stacks-chow/tags.txt, source lines 3807--3807. Prerequisite pointer invoked at proof line(s) 1005; it names a standard fact of a sibling chapter that the counted argument uses, and every step of the counted proof built on it is retained verbatim in the delivered body.\par\medskip
\ifdefined\hypertarget\hypertarget{sourcepointer-6}{}\fi
\noindent\textit{Source reference.} Original lemma 02X4: sources/stacks-chow/tags.txt, source lines 3827--3827. Prerequisite pointer invoked at proof line(s) 1022; it names a standard fact of a sibling chapter that the counted argument uses, and every step of the counted proof built on it is retained verbatim in the delivered body.\par\medskip
\ifdefined\hypertarget\hypertarget{sourcepointer-7}{}\fi
\noindent\textit{Source reference.} Original lemma 02WZ: sources/stacks-chow/tags.txt, source lines 3822--3822. Prerequisite pointer invoked at proof line(s) 1034; it names a standard fact of a sibling chapter that the counted argument uses, and every step of the counted proof built on it is retained verbatim in the delivered body.\par\medskip
\ifdefined\hypertarget\hypertarget{sourcepointer-8}{}\fi
\noindent\textit{Source reference.} Original lemma 02YO: sources/stacks-chow/tags.txt, source lines 3883--3883. Prerequisite pointer invoked at proof line(s) 1008; it names a standard fact of a sibling chapter that the counted argument uses, and every step of the counted proof built on it is retained verbatim in the delivered body.\par\medskip
\begin{lemma}
\label{lemma-separated-cover}
Let $S$ be a scheme.
Let $X$ be an algebraic space over $S$.
Let $U$ be a separated scheme and $U \to X$ \'etale.
Then $U \to X$ is separated, and $R = U \times_X U$ is a separated scheme.
\end{lemma}

\begin{proof}
Let $X' \subset X$ be the open subscheme such that $U \to X$ factors
through an \'etale surjection $U \to X'$, see
Lemma \ref{lemma-etale-image-open}.
If $U \to X'$ is separated, then so is $U \to X$, see
Spaces, Lemma
\ref{spaces-lemma-composition-representable-transformations-property}
(as the open immersion $X' \to X$ is separated by
Spaces, Lemma
\ref{spaces-lemma-representable-transformations-property-implication}
and
Schemes, Lemma \ref{schemes-lemma-immersions-monomorphisms}).
Moreover, since $U \times_{X'} U = U \times_X U$ it suffices
to prove the result after replacing $X$ by $X'$, i.e., we may
assume $U \to X$ surjective.
Consider the commutative diagram
$$
\xymatrix{
R = U \times_X U \ar[r] \ar[d] & U \ar[d] \\
U \ar[r] & X
}
$$
In the proof of
Spaces, Lemma \ref{spaces-lemma-properties-diagonal}
we have seen that $j : R \to U \times_S U$ is separated.
The morphism of schemes $U \to S$ is separated as $U$ is a separated
scheme, see
Schemes, Lemma \ref{schemes-lemma-compose-after-separated}.
Hence $U \times_S U \to U$ is separated as a base change, see
Schemes, Lemma \ref{schemes-lemma-separated-permanence}.
Hence the scheme $U \times_S U$ is separated (by the same lemma).
Since $j$ is separated we see in the same way that $R$ is separated.
Hence $R \to U$ is a separated morphism (by
Schemes, Lemma \ref{schemes-lemma-compose-after-separated}
again). Thus by
Spaces, Lemma \ref{spaces-lemma-representable-morphisms-spaces-property}
and the diagram above we conclude that $U \to X$ is separated.
\end{proof}

\end{document}
